\begin{trapbox}
\begin{center}\begin{tikzpicture}[>=stealth,font=\small]
\node[anchor=west,font=\small\bfseries] at (-.6,4.95) {同一公共点：水平线与真正的切线};
\draw[->,gray] (-.6,0)--(4.2,0) node[right] {$x$};
\draw[->,gray] (0,-1.2)--(0,4.1) node[above] {$y$};
\draw[blue,very thick,domain=-1.1:3.9,samples=100] plot ({\x*\x/4},\x);
\node[blue,fill=white] at (3.3,3.05) {$y^2=4x$};
\node[below left] at (0,0) {$O$};
\draw[red,dashed,thick] (-.5,2)--(4,2) node[right] {$y=2$：非切线};
\draw[orange!80!black,very thick] (-.5,.5)--(3,4) node[right] {$y=x+1$};
\draw[orange!80!black,dotted,thick] (0,-1.1)--(0,4);
\fill (1,2) circle (2pt) node[below right=7pt] {$Q(1,2)$};
\node[anchor=west] at (5,3) {两条直线都只交于 $Q$};
\node[anchor=west] at (5,1) {只有橙色实线在 $Q$ 处相切};
\node[anchor=west] at (5,-.1) {橙色点线：顶点切线 $x=0$};
\end{tikzpicture}
\end{center}
    \begin{description}[style=nextline, leftmargin=2.8em, labelsep=0.5em, itemsep=0.38em]
        \item[\textbf{\textcolor{CTrap}{易错点一（把过点当切点）}}]
        题目说“切线经过$A(4,5)$”，就直接令$x_0=4,y_0=5$。
        \item[\textbf{\textcolor{CTrap}{正确理解：}}]
        切点式里的$P(x_0,y_0)$必须在抛物线上。比如对$y^2=4x$，
        $5^2\ne4\times4$，所以$A(4,5)$不是切点；应另设切点并联立条件求解。
    \end{description}

    \medskip
    \begin{description}[style=nextline, leftmargin=2.8em, labelsep=0.5em, itemsep=0.38em]
        \item[\textbf{\textcolor{CTrap}{易错点二（记错开口方向）}}]
        把$y^2=2px$的公式直接用到$x^2=2py$或$y^2=-2px$。
        \item[\textbf{\textcolor{CTrap}{正确理解：}}]
        横开口时对应$y_0y=\pm p(x+x_0)$，竖开口时对应$x_0x=\pm p(y+y_0)$；
        $\pm$与标准方程同号。四种情况都需检验切点在对应曲线上。
    \end{description}

    \medskip
    \begin{description}[style=nextline, leftmargin=2.8em, labelsep=0.5em, itemsep=0.38em]
        \item[\textbf{\textcolor{CTrap}{易错点三（顶点除以零）}}]
        用$k=p/y_0$代入$y_0=0$，误判顶点切线不存在。
        \item[\textbf{\textcolor{CTrap}{正确理解：}}]
        $y_0=0$时$x_0=0$，切点式给出$x=0$。\textbf{切线存在，但斜率不存在。}
        若采用导数法，对$x=\dfrac{y^2}{2p}$求导，不会漏掉竖直切线。
    \end{description}

    \medskip
    \begin{description}[style=nextline, leftmargin=2.8em, labelsep=0.5em, itemsep=0.38em]
        \item[\textbf{\textcolor{CTrap}{易错点四（一个交点就等于相切）}}]
        认为“联立后只有一组坐标”就证明直线相切。
        \item[\textbf{\textcolor{CTrap}{最小反例：}}]
        抛物线$y^2=4x$与水平线$y=2$仅有公共点$(1,2)$，
        但在该点的切线是$y=x+1$，不是$y=2$。
        \item[\textbf{\textcolor{CTrap}{正确理解：}}]
        通用证明中先设$x=my+b$，消去$x$后得到\textbf{关于$y$的一元二次方程}，
        这时用$\Delta=0$判断相切；不能把$y=c$代入得到的\textbf{一次方程只有一个解}误当成判别式为零。
    \end{description}

    \medskip
    \begin{description}[style=nextline, leftmargin=2.8em, labelsep=0.5em, itemsep=0.38em]
        \item[\textbf{\textcolor{CTrap}{易错点五（设直线漏掉顶点）}}]
        只设$y=kx+b$，因斜率不存在而遗漏竖直直线$x=0$。
        \item[\textbf{\textcolor{CTrap}{正确理解：}}]
        在$y^2=2px$中使用$x=my+b$可统一处理顶点；或直接使用切点式。
    \end{description}


    \medskip
    \begin{description}[style=nextline, leftmargin=2.8em, labelsep=0.5em, itemsep=0.38em]
        \item[\textbf{\textcolor{CTrap}{易错点六（把截距误写成$x_0$）}}]
        认为切线与$x$轴交点是$(x_0,0)$，把切点在$x$轴上的投影$H$误当成截点$T$。
        \item[\textbf{\textcolor{CTrap}{正确理解：}}]
        $H(x_0,0)$是切点的投影，$T(-x_0,0)$才是切线与$x$轴的交点。
        对非顶点切点，$T,O,H$依次在$x$轴上，且$TO=OH=x_0$。
    \end{description}

    \medskip
    \begin{description}[style=nextline, leftmargin=2.8em, labelsep=0.5em, itemsep=0.38em]
        \item[\textbf{\textcolor{CTrap}{易错点七（忽略顶点与开口方向）}}]
        把$P=O$时的$T=H=O$当成长度大于零的线段，或认为任意抛物线的截距都必须在$x$轴。
        \item[\textbf{\textcolor{CTrap}{正确理解：}}]
        顶点处仍有$x$轴截距$0=-x_0$，但$T,H,O$重合。
        对$y^2=\pm2px$，相关截距在$x$轴上，等于$-x_0$；
        对$x^2=\pm2py$，应当看\textbf{$y$轴截距}，等于$-y_0$。
    \end{description}
\end{trapbox}
